 

%请使用xelatex编译
\documentclass[11pt,addpoints,answers]{exam}


% !Mode:: "TeX:UTF-8"
% !TEX root=./hw1.tex
\usepackage{xeCJK}    % 写中文要用到
\usepackage{tikz}
\usepackage{amsmath,amsfonts,amssymb,amsthm}
\usepackage{fullpage}
\usepackage{times}
\usepackage{hyperref}
\usepackage{pdfsync}
\usepackage{microtype}
\usepackage{color}
\usepackage{cleveref}
\crefformat{footnote}{#2\footnotemark[#1]#3}
\definecolor{light-gray}{gray}{0.5}

\renewcommand{\tablename}{表}
\renewcommand{\figurename}{图}
\renewcommand{\proof}{证明}

% \theoremstyle{plain}
% \newtheorem{theorem}{定理~}
% \newtheorem{lemma}{引理~}
% \newtheorem{axiom}{公理~}
% \newtheorem{proposition}{命题~}
% \newtheorem{prop}{性质~}
% \newtheorem{corollary}{推论~}
% \newtheorem{conclusion}{结论~}
% \newtheorem{definition}{定义~}
% \newtheorem{conjecture}{猜想~}
% \newtheorem{example}{例~}
% \newtheorem{remark}{注~}
% \newtheorem{algorithm}{算法~}


%%% BLACKBOARD SYMBOLS

% \newcommand{\C}{\ensuremath{\mathbb{C}}}
\newcommand{\D}{\ensuremath{\mathbb{D}}}
\newcommand{\F}{\ensuremath{\mathbb{F}}}
% \newcommand{\G}{\ensuremath{\mathbb{G}}}
\newcommand{\J}{\ensuremath{\mathbb{J}}}
\newcommand{\N}{\ensuremath{\mathbb{N}}}
\newcommand{\Q}{\ensuremath{\mathbb{Q}}}
\newcommand{\R}{\ensuremath{\mathbb{R}}}
\newcommand{\T}{\ensuremath{\mathbb{T}}}
\newcommand{\Z}{\ensuremath{\mathbb{Z}}}
\newcommand{\QR}{\ensuremath{\mathbb{QR}}}

\newcommand{\Zt}{\ensuremath{\Z_t}}
\newcommand{\Zp}{\ensuremath{\Z_p}}
\newcommand{\Zq}{\ensuremath{\Z_q}}
\newcommand{\ZN}{\ensuremath{\Z_N}}
\newcommand{\Zps}{\ensuremath{\Z_p^*}}
\newcommand{\ZNs}{\ensuremath{\Z_N^*}}
\newcommand{\JN}{\ensuremath{\J_N}}
\newcommand{\QRN}{\ensuremath{\QR_{N}}}
\newcommand{\QRp}{\ensuremath{\QR_{p}}}

%%% THEOREM COMMANDS

\theoremstyle{plain}            % following are "theorem" style
\newtheorem{theorem}{定理}[section]
\newtheorem{lemma}[theorem]{引理}
\newtheorem{corollary}[theorem]{推论}
\newtheorem{proposition}[theorem]{命题}
\newtheorem{claim}[theorem]{声明}
\newtheorem{fact}[theorem]{事实}

\theoremstyle{definition}       % following are def style
\newtheorem{definition}[theorem]{定义}
\newtheorem{conjecture}[theorem]{推测}
\newtheorem{example}[theorem]{例}
\newtheorem{protocol}[theorem]{协议}

\theoremstyle{remark}           % following are remark style
\newtheorem{remark}[theorem]{注}
\newtheorem{note}[theorem]{注}
\newtheorem{exercise}[theorem]{练习}

% equation numbering style
\numberwithin{equation}{section}

%%% GENERAL COMPUTING

\newcommand{\bit}{\ensuremath{\set{0,1}}}
\newcommand{\pmone}{\ensuremath{\set{-1,1}}}

% asymptotics
\DeclareMathOperator{\poly}{poly}
\DeclareMathOperator{\polylog}{polylog}
\DeclareMathOperator{\negl}{negl}
\newcommand{\Otil}{\ensuremath{\tilde{O}}}

% probability/distribution stuff
\DeclareMathOperator*{\E}{E}
\DeclareMathOperator*{\Var}{Var}

% sets in calligraphic type
\newcommand{\calA}{\ensuremath{\mathcal{A}}}
\newcommand{\calD}{\ensuremath{\mathcal{D}}}
\newcommand{\calF}{\ensuremath{\mathcal{F}}}
\newcommand{\calH}{\ensuremath{\mathcal{H}}}
\newcommand{\calK}{\ensuremath{\mathcal{K}}}
\newcommand{\calM}{\ensuremath{\mathcal{M}}}
\newcommand{\calX}{\ensuremath{\mathcal{X}}}
\newcommand{\calY}{\ensuremath{\mathcal{Y}}}

% types of indistinguishability
\newcommand{\compind}{\ensuremath{\stackrel{c}{\approx}}}
\newcommand{\statind}{\ensuremath{\stackrel{s}{\approx}}}
\newcommand{\perfind}{\ensuremath{\equiv}}

% font for general-purpose algorithms
\newcommand{\algo}[1]{\ensuremath{\mathsf{#1}}}
% font for general-purpose computational problems
\newcommand{\problem}[1]{\ensuremath{\mathsf{#1}}}
% font for complexity classes
\newcommand{\class}[1]{\ensuremath{\mathsf{#1}}}

% complexity classes and languages
\renewcommand{\P}{\class{P}}
\newcommand{\BPP}{\class{BPP}}
\newcommand{\NP}{\class{NP}}
\newcommand{\coNP}{\class{coNP}}
\newcommand{\AM}{\class{AM}}
\newcommand{\coAM}{\class{coAM}}
\newcommand{\IP}{\class{IP}}

%%% "LEFT-RIGHT" PAIRS OF SYMBOLS

% inner product
\newcommand{\inner}[1]{\langle{#1}\rangle}
\newcommand{\innerfit}[1]{\left\langle{#1}\right\rangle}
% absolute value
\newcommand{\abs}[1]{\lvert{#1}\rvert}
\newcommand{\absfit}[1]{\left\lvert{#1}\right\rvert}
% a set
\newcommand{\set}[1]{\{{#1}\}}
\newcommand{\setfit}[1]{\left\{{#1}\right\}}
% parens
\newcommand{\parens}[1]{({#1})}
\newcommand{\parensfit}[1]{\left({#1}\right)}
% tuple = alias for parens
\newcommand{\tuple}[1]{\parens{#1}}
\newcommand{\tuplefit}[1]{\parensfit{#1}}
% square brackets
\newcommand{\bracks}[1]{[{#1}]}
\newcommand{\bracksfit}[1]{\left[{#1}\right]}
% rounding off
\newcommand{\round}[1]{\lfloor{#1}\rceil}
% floor function
\newcommand{\floor}[1]{\lfloor{#1}\rfloor}
% ceiling function
\newcommand{\ceil}[1]{\lceil{#1}\rceil}
% length of a string
\newcommand{\len}[1]{\lvert{#1}\rvert}
\newcommand{\lenfit}[1]{\left\lvert{#1}\right\rvert}
% length of some vector, element
\newcommand{\length}[1]{\lVert{#1}\rVert}
\newcommand{\lengthfit}[1]{\left\lVert{#1}\right\rVert}

%%% CRYPTO-RELATED NOTATION

% KEYS AND RELATED

\newcommand{\key}[1]{\ensuremath{#1}}

\newcommand{\pk}{\key{pk}}
\newcommand{\vk}{\key{vk}}
\newcommand{\sk}{\key{sk}}
\newcommand{\mpk}{\key{mpk}}
\newcommand{\msk}{\key{msk}}
\newcommand{\fk}{\key{fk}}
\newcommand{\id}{id}
\newcommand{\keyspace}{\ensuremath{\mathcal{K}}}
\newcommand{\msgspace}{\ensuremath{\mathcal{M}}}
\newcommand{\ctspace}{\ensuremath{\mathcal{C}}}
\newcommand{\tagspace}{\ensuremath{\mathcal{T}}}
\newcommand{\idspace}{\ensuremath{\mathcal{ID}}}

\newcommand{\concat}{\ensuremath{\|}}

% GAMES

% advantage
\newcommand{\advan}{\ensuremath{\mathbf{Adv}}}

% different attack models
\newcommand{\attack}[1]{\ensuremath{\text{#1}}}

\newcommand{\atk}{\attack{atk}} % dummy attack
\newcommand{\indcpa}{\attack{ind-cpa}}
\newcommand{\indcca}{\attack{ind-cca}}
\newcommand{\anocpa}{\attack{ano-cpa}} % anonymous
\newcommand{\anocca}{\attack{ano-cca}}
\newcommand{\euacma}{\attack{eu-acma}} % forgery: adaptive chosen-message
\newcommand{\euscma}{\attack{eu-scma}} % forgery: static chosen-message
\newcommand{\suacma}{\attack{su-acma}} % strongly unforgeable

% ADVERSARIES
\newcommand{\attacker}[1]{\ensuremath{\mathcal{#1}}}

\newcommand{\Adv}{\attacker{A}}
\newcommand{\AdvA}{\attacker{A}}
\newcommand{\AdvB}{\attacker{B}}
\newcommand{\Dist}{\attacker{D}}
\newcommand{\Sim}{\attacker{S}}
\newcommand{\Ora}{\attacker{O}}
\newcommand{\Inv}{\attacker{I}}
\newcommand{\For}{\attacker{F}}

% CRYPTO SCHEMES

\newcommand{\scheme}[1]{\ensuremath{\text{#1}}}

% pseudorandom stuff
\newcommand{\prg}{\algo{PRG}}
\newcommand{\prf}{\algo{PRF}}
\newcommand{\prp}{\algo{PRP}}

% symmetric-key cryptosystem
\newcommand{\skc}{\scheme{SKC}}
\newcommand{\skcgen}{\algo{Gen}}
\newcommand{\skcenc}{\algo{Enc}}
\newcommand{\skcdec}{\algo{Dec}}

% public-key cryptosystem
\newcommand{\pkc}{\scheme{PKC}}
\newcommand{\pkcgen}{\algo{Gen}}
\newcommand{\pkcenc}{\algo{Enc}} % can also use \kemenc and \kemdec
\newcommand{\pkcdec}{\algo{Dec}}

% digital signatures
\newcommand{\sig}{\scheme{SIG}}
\newcommand{\siggen}{\algo{Gen}}
\newcommand{\sigsign}{\algo{Sign}}
\newcommand{\sigver}{\algo{Ver}}

% message authentication code
\newcommand{\mac}{\scheme{MAC}}
\newcommand{\macgen}{\algo{Gen}}
\newcommand{\mactag}{\algo{Tag}}
\newcommand{\macver}{\algo{Ver}}

% key-encapsulation mechanism
\newcommand{\kem}{\scheme{KEM}}
\newcommand{\kemgen}{\algo{Gen}}
\newcommand{\kemenc}{\algo{Encaps}}
\newcommand{\kemdec}{\algo{Decaps}}

% identity-based encryption
\newcommand{\ibe}{\scheme{IBE}}
\newcommand{\ibesetup}{\algo{Setup}}
\newcommand{\ibeext}{\algo{Ext}}
\newcommand{\ibeenc}{\algo{Enc}}
\newcommand{\ibedec}{\algo{Dec}}

% hierarchical IBE (as key encapsulation)
\newcommand{\hibe}{\scheme{HIBE}}
\newcommand{\hibesetup}{\algo{Setup}}
\newcommand{\hibeext}{\algo{Extract}}
\newcommand{\hibeenc}{\algo{Encaps}}
\newcommand{\hibedec}{\algo{Decaps}}

% binary tree encryption (as key encapsulation)
\newcommand{\bte}{\scheme{BTE}}
\newcommand{\btesetup}{\algo{Setup}}
\newcommand{\bteext}{\algo{Extract}}
\newcommand{\bteenc}{\algo{Encaps}}
\newcommand{\btedec}{\algo{Decaps}}

% trapdoor functions
\newcommand{\tdf}{\scheme{TDF}}
\newcommand{\tdfgen}{\algo{Gen}}
\newcommand{\tdfeval}{\algo{Eval}}
\newcommand{\tdfinv}{\algo{Invert}}
\newcommand{\tdfver}{\algo{Ver}}

%%% PROTOCOLS

\newcommand{\out}{\text{out}}
\newcommand{\view}{\text{view}}

%%% COMMANDS FOR LECTURES/HOMEWORKS

\newcommand{\lecheader}{%
  \chead{\large \textbf{Lecture \lecturenum\\\lecturetopic}}

  \lhead{\small
    \textbf{\href{http://bhpan.buaa.edu.cn}{研究生课程：计算复杂性}\\2023年秋季}}

  \rhead{\small \textbf{主讲: 张宗洋
      }

  \setlength{\headheight}{20pt}
  \setlength{\headsep}{16pt}
}}

\newcommand{\hwheader}{%
  \chead{\Large \textbf{作业 \hwnum}}

  \lhead{\small
    \textbf{{计算复杂性}\\2025年秋季}}

  \rhead{\small \textbf{主讲: 张宗洋
      \\姓名: \studentname}}

  \setlength{\headheight}{20pt}
  \setlength{\headsep}{16pt}
  
  \headrule
}

\newcommand{\examheader}{%
  \chead{\Large \textbf{测试 \\ \duedate}}
  
\lhead{\small
    \textbf{\href{http://bhpan.buaa.edu.cn}{计算复杂性}\\2024年秋季}}

  \rhead{\small \textbf{主讲: 张宗洋
      \\姓名: \studentname}}

  \setlength{\headheight}{20pt}
  \setlength{\headsep}{16pt}
  
  \headrule
}
  
	
\newcommand{\hint}[1]{\href{http://www.cims.nyu.edu/~regev/cgi-bin/hints/index.html?#1}{{\textcolor{light-gray}{I need a hint! (ID #1)}}}}

\newcommand{\hinttext}[2]{\href{http://www.cims.nyu.edu/~regev/cgi-bin/hints/index.html?#1}{{\textcolor{light-gray}{#2 (ID #1)}}}}

\newcommand{\hintcost}[2]{\href{http://www.cims.nyu.edu/~regev/cgi-bin/hints/index.html?#1}{{\textcolor{light-gray}{I need a hint for #2 points! (ID #1)}}}}
	
\newcommand{\moreinfo}[1]{\href{http://www.cims.nyu.edu/~regev/cgi-bin/hints/index.html?#1}{{\textcolor{light-gray}{I'm done solving and want to know more! (ID #1)}}}}



\newcommand{\hwnum}{3}
\newcommand{\duedate}{11月24日} % changing this does  e :)
\newcommand{\studentname}{}

% END OF SUPPLIED VARIABLES

\hwheader                       % execute homework commands

\begin{document}
	
 
	
	\pagestyle{head}                % put header on every page
\begin{itemize}
	\item 请将PDF文件发送至邮箱：complexity\_buaa@163.com
	\item 文件命名规则是``学号-姓名-第三次作业.pdf"
	\item 答案请用\LaTeX，打分基于认真性、正确性和描述的清晰性
	\item 务必独立自主解题
	\item 作业截止时间为 \textbf{\duedate}.
\end{itemize}


	\begin{questions}
		
		
		\question P118页, 3.2题。此练习与图灵机$M_1$有关, 例3.5给出了它的描述及状态图。在下列每个输入串上，给出$M_1$所进入的格局序列。
		
		\begin{parts}
			
			\part (c)$1\#\#1$
			\begin{solution}
				
                $q_11\#\#1$
                
                $xq_3\#\#1$
                
                $x\#q_5\#1$

                $x\#\#q_r$
                
			\end{solution}
			
			\part (e)$10\#10$
			\begin{solution}

                $q_110\#10$

                $xq_30\#10$

                $x0q_3\#10$

                $x0\#q_510$

                $x0q_6\#x0$

                $xq_70\#x0$

                $q_7x0\#x0$

                $xq_10\#x0$

                $xxq_2\#x0$

                $xx\#q_4x0$

                $xx\#xq_40$

                $xx\#q_6xx$

                $xxq_6\#xx$

                $xq_7x\#xx$

                $xxq_1\#xx$

                $xx\#q_8xx$

                $xx\#xq_8x$

                $xx\#xxq_8$

                $xx\#xxq_{acc}$
			\end{solution}
			
			
			
		\end{parts}
		
		\question P119页, 3.8.下面的语言都是字母表$\{0,1\}$上的语言, 以实现层次的描述给出判定这些语言的图灵机：
		\begin{parts}	
			\part	(b). $\{\omega|\omega\text{所包含的0的个数是1的个数的两倍}\}$
			\begin{solution}
				
                设计双带图灵机：
                
                第一条带子为输入序列。
                第二条带子读写头不动，只用第一个位置记录0个数的两倍和1的个数之差。
                
                初始将第二条带子第一个位置写为0.
                
                从左向右扫描第一条带子，如果当前位置输入是0，则第二条带子上第一个位置数加2；如果为1，则减1。
                
                当读取到输入终止符时，检查第二个带子上的计数，如果是0则接受，反之不接受。
			\end{solution}
			
			
			\part	(c). $\{\omega|\omega\text{所包含的0的个数不是1的个数的两倍}\}$
			\begin{solution}
				使用第一题能够判断“0的个数是1的个数的两倍”的图灵机 M。
                若M接受，则拒绝；若M拒绝，则接受。
			\end{solution}
		\end{parts}
		
		
		\question P119页, 3.11题.
		证明： 不能在带子的输入区域写的单带图灵机只能识别正则语言。
		\begin{solution}
			
            若图灵机不能在输入区域写，则它相当于一个只读图灵机。            
            
            只读图灵机的行为等价于有限自动机，其状态转移仅依赖于当前读头位置和状态。
            
            因此，它只能识别正则语言。
			
		\end{solution}  
		
		
		\question P119页, 3.12题.
		证明每一个无穷图灵可识别语言都有一个无穷可判定子集。
		\begin{solution}
			
            设 L 是无穷的图灵可识别语言，由枚举器 E 枚举。
            
            构造子集 S：选取 E 输出的第一个串 w1，然后选取下一个比 w1 长的串 w2，依此类推。
            
            S 是无穷的，且是可判定的，因为我们可以判断一个串是否在 S 中（模拟 E 直到输出该串或更长的串）。
			
		\end{solution}  
		
		\question P119页, 3.14题.
		设图灵可识别语言$B=\left\{\langle M_{1}\rangle ,\langle M_{2}\rangle ... \right\}$由图灵机的描述组成，证明存在一个以图灵机作为输入字符的图灵可识别语言$C$，使得在$B$中的每一个图灵机在$C$中也有一个图灵机与之等价，反之亦然。
		\begin{solution}
			构造 C = $\{$ ⟨M⟩ | M 等价于 B 中某个 Mi $\}$。
            
            可以构造一个枚举器，枚举所有与 B 中机器等价的图灵机。
            
            因此 C 是图灵可识别的。
			
		\end{solution}  
		
		\question P119页, 3.15题
		证明图灵可识别语言类在下列运算下封闭：
		a.并 b.连接 c.星号 d.交 e.同态
		\begin{solution}
			a. 并：使用双带图灵机模拟两个识别器，若有一个接受则接受。  
            
            b. 连接：非确定性地猜测分割点，模拟两个识别器。  
            
            c. 星号：非确定性地猜测分割点并模拟多次。  
            
            d. 交：使用双带图灵机模拟两个识别器，若都接受则接受。  
            
            e. 同态：对输入应用同态后模拟原识别器。
			
		\end{solution}  
		
		\question P119页, 3.16题
		证明可判定语言类在下列运算下封闭：
		a.并 b.连接 c.星号 d.补 e.交
		\begin{solution}
			
            a. 并：模拟两个判定器，若有一个接受则接受。  

            b. 连接：枚举所有分割点，模拟两个判定器。  
            
            c. 星号：动态规划判断是否可由多个子串连接而成。  
            
            d. 补：交换接受和拒绝状态。  
            
            e. 交：模拟两个判定器，若都接受则接受。
			
		\end{solution}   
		
		
		\question P119页, 3.17题
		只写一次图灵机 （write-once Turing machine) 是一个单带图灵机， 它在每个带子方格上最多只
		能改变其内容一次（包括带子上的输入区）。 证明图灵机模型的这个变形等价于普通的图灵机模
		型。 （ 提示： 首先考虑如下图灵机： 可以修改带子方格最多两次， 使用多个带子。）
		\begin{solution}
			
            使用带子记录修改历史。
            
            每次“写一次”限制被违反时，将所有内容复制到带子后端并用$\#$和原先内容隔开继续模拟。
            
            因此只写一次图灵机可以模拟普通图灵机。
			
		\end{solution}  
		
		\question P119页, 3.18题
		双无限带图灵机 （Turing machine with doubly infinite tape) 与普通图灵机相似， 所不同的是它的
		带子向左和向右都是无限的。 此带子在开始时， 除了包括输人区域外 ， 其他都填以空白符， 计
		算也像通常一样定义， 只是在它向左移动时不会遇到带子的端点。 证明这种类型的图灵机识别
		图灵可识别语言类。
		\begin{solution}
			
            用普通图灵机模拟双无限带：将双无限带对折成两个半带，用两个带子模拟。
            
            反之，双无限带图灵机只需要记录读写头初始位置，就可以模拟普通图灵机。
            
            因此两者等价。
			
		\end{solution} 
		
			\question P119页, 3.19题。
		左复位图灵机与普通图灵机相似， 只是它的转移函数具有以下形式：
		\[
		\delta:Q\times T\rightarrow Q\times T\times \{R,\text{RESET}\}
		\]
		如果 $\delta(q,a)=(r,b,\text{RESET})$，则当机器处于状态 $q$ 且读 $a$ 时，在带上上写下 $b$ 并进入状态 $r$后，读写头就跳到带子的左端点。证明左复位图灵机识别图灵可识别语言类。
		\begin{solution}
			用普通图灵机模拟左复位图灵机：在带上标记复位位置，每次复位时移动到头。
            
            反之，左复位图灵机只需要忽略 RESET 功能，就可以模拟普通图灵机。
            
            因此两者等价。
			
		\end{solution}
	\end{questions}
	
\end{document}
